\input{estilo}
% No está en la pizarra: dibuja el diente de sierra que el texto describe con
% los números del paper (§3): get time cada 30 s, 200 µs/s de drift aplicado,
% ε de 1 a 7 ms y casi siempre ≈ 4. A escala: x = 0,06 cm por segundo,
% y = 0,4 cm por milisegundo.
\begin{document}
\begin{figura}\begin{tikzpicture}[x=0.06cm, y=0.4cm]
  \draw[->] (0,0) -- (100,0) node[etiq, right] {tiempo (\unid{s})};
  \draw[->] (0,0) -- (0,8.6) node[etiq, above] {incertidumbre \var{ε} (\unid{ms})};
  % cada get time, en el eje: sin flechas que lo crucen
  \draw (0,0) -- (0,-0.2) node[etiq, below] {0};
  \foreach \t in {30,60,90}
    \draw (\t,0) -- (\t,-0.2) node[etiq, below, align=center] {\t\\[1pt]{\color{segundo}get time}};
  \foreach \e in {1,7} \draw (0,\e) -- (-1.5,\e) node[etiq, left] {\e};
  \draw (0,4) -- (-1.5,4) node[etiq, left] {{\color{suave}casi siempre ≈}\ 4};
  % ε: crece 6 ms en 30 s y cae al pedir de nuevo
  \draw[segundo, line width=1pt] (0,1) -- (30,7) -- (30,1) -- (60,7) -- (60,1) -- (90,7) -- (90,1) -- (94,1.8);
  % la cuenta, a la derecha de todo
  \draw[llave=4pt] (104,7) -- (104,1);
  \node[etiq, align=left, text=suave, anchor=west] at (107,4) {\unid{30 s × 200 µs/s}\\\unid{= 6 ms}};
\end{tikzpicture}\end{figura}
\end{document}
